\subsection{Конституция Российской Федерации}

Конституция Российской Федерации является основным нормативным правовым актом, обладающим высшей юридической силой. Именно её положения формируют базовые принципы обращения с информацией, обеспечением неприкосновенности частной жизни и защитой прав граждан в информационной сфере. В частности, Конституция закрепляет право каждого на неприкосновенность частной жизни, личную и семейную тайну, защиту своей чести и доброго имени, а также устанавливает свободу поиска, получения, передачи, производства и распространения информации любым законным способом.

В контексте обработки audit-логов Kubernetes-кластера особенно актуальны следующие статьи Конституции РФ:

\begin{itemize}
    \item \textbf{Статья 23} — закрепляет право каждого на неприкосновенность частной жизни, личную и семейную тайну, защиту своей чести и доброго имени. Это означает, что любые данные, которые позволяют идентифицировать действия конкретного пользователя или служебной учётной записи, должны обрабатываться с соблюдением принципа конфиденциальности.
    \item \textbf{Статья 24} — устанавливает, что сбор, хранение, использование и распространение информации о частной жизни лица без его согласия не допускаются. В контексте разработки системы обнаружения аномалий это накладывает требования к обезличиванию или псевдонимизации персональных данных, а также к ограничению доступа к audit-логам только уполномоченным сотрудникам.
\end{itemize}

Таким образом, при проектировании и внедрении системы обнаружения аномалий необходимо учитывать конституционные принципы, обеспечивающие защиту персональной и служебной информации. Любые действия по анализу и обработке логов должны быть законными, прозрачными и документированными, с соблюдением требований по минимизации сбора и обработки личных данных. Реализация программного средства обнаружения аномалий в Kubernetes-кластере должна строиться с учетом следующих конституционных принципов:

\begin{enumerate}
    \item Обеспечение конфиденциальности персональных данных пользователей и служебной информации.
    \item Законность и прозрачность методов сбора и обработки данных.
    \item Ограничение доступа к информации только уполномоченным лицам.
    \item Документирование всех процедур обработки логов с целью соблюдения права граждан на защиту своей частной жизни.
\end{enumerate}